\documentclass[10pt]{article}
\usepackage{icmlko}

\icmlmeta{IP 파운데이션 월드모델}{270}{시계 몫}{쇠퇴가 제일 큰 도메인과 시계에 지배당한 도메인이 다르다 --- 게임은 34배 쇠퇴하는데 멀쩡하다}

\begin{document}

\twocolumn[
\icmlheader

\begin{abstract}
노트 269가 도서 라벨이 시계임을 밝히고 ``같은 검사를 아홉 곳에''를 남겼다.
모형을 안 쓰는 검사라 값싸다 --- 유보를 반년 코호트로 갈라 라벨 중앙값의
첫 조각 대 끝 조각 비를 잰다. \textbf{제일 크게 쇠퇴하는 것은 도서가 아니라
게임이다}(34.5배 대 10.5배). 그런데 게임은 멀쩡하다. 왜냐면 \textbf{쇠퇴의
크기와 순위가 날짜에 붙는 정도는 다른 것}이기 때문이다 --- 게임 라벨은
노트 210이 갈아 끼운 고정 30일 창(\texttt{y\_w30})이라 최근 게임의 리뷰가
적기는 해도 \emph{순위}가 날짜에 안 붙는다(날짜 상관 $+$0.182). 도서는
쇠퇴가 그 3분의 1인데 날짜 상관이 $+$0.562 다. \textbf{쓸 눈금은 쇠퇴비가
아니라 달력 대 모형의 비}이고, 그것을 \texttt{rank.clock} 으로 박았다.
\textbf{여덟 중 도서 하나만 1 을 넘는다}(1.52) --- 즉 도서 점수는 능력이
아니라 시계다. 웹툰이 0.77 로 둘째인데 1 아래라 모형이 이긴다. 나머지 여섯은
0.37 이하이고 팝업 · 펀딩은 음수다(오래된 것이 오히려 낮다). 부수적으로
게임 라벨의 잘림 처리가 옳은 것도 확인했다 --- 30일이 안 찬 43건이 전부
결측으로 빠져 있다. 그리고 \texttt{rank.clock} 을 처음 돌렸을 때 \textbf{게임이
표에서 빠졌다} --- ``유보 행 수''와 ``라벨이 있는 행 수''를 같은 것으로 보고
길이를 견줬기 때문이다. 그 43건이 정확히 그 차이였다.
\end{abstract}
]

\section{값싼 검사를 아홉 곳에}

노트 269의 검사는 모형을 안 쓴다 --- 유보를 반년 코호트로 갈라 라벨 중앙값이
첫 조각에서 끝 조각으로 얼마나 떨어지나 본다. 아홉에 다 돌렸다(만화 6건 ·
아이돌 25건은 표본 부족).

\begin{center}\small
\begin{tabular}{lrrr}
\hline
도메인 & 유보 & 쇠퇴비 & $-$날짜 상관 \\
\hline
\textbf{게임} & 180 & \textbf{34.51} & $+$0.182 \\
모바일 & 441 & 12.69 & $+$0.198 \\
\textbf{도서} & 163 & 10.46 & $\mathbf{+0.562}$ \\
웹툰 & 711 & 4.78 & $+$0.350 \\
애니 & 606 & 1.62 & $+$0.144 \\
세계애니 & 300 & 1.17 & $+$0.099 \\
펀딩 & 80 & 0.92 & $-$0.012 \\
팝업 & 59 & 0.52 & $-$0.113 \\
\hline
\end{tabular}
\end{center}

\textbf{제일 크게 쇠퇴하는 것은 게임이다} --- 유보 안에서 34.5배. 도서(10.5배)의
세 배가 넘는다. 그런데 노트 267 · 269가 문제라고 짚은 것은 도서였다.

\section{쇠퇴가 크다고 시계인 건 아니다}

두 수가 다른 것을 재고 있다.

\begin{quote}\small
\textbf{쇠퇴비}는 ``최근 것의 라벨이 얼마나 작나''이고,
\textbf{날짜 상관}은 ``라벨 \emph{순위}가 날짜를 얼마나 따르나''다.
\end{quote}

게임 라벨은 노트 210이 누적 리뷰에서 갈아 끼운 \textbf{고정 30일 창}
(\texttt{y\_w30})이다. 최근 게임이 30일에 받는 리뷰가 옛 게임보다 적기는
하지만 --- 스팀 신작 수가 늘어 한 작품이 받는 몫이 준다 --- 그 크기 차이가
\emph{순위}를 날짜로 정렬하지는 않는다. 그래서 쇠퇴비 34.5 에 날짜 상관
$+$0.182 다.

도서는 반대다. \texttt{sales\_point} 는 출간부터 계속 쌓이므로 \textbf{늦게
나올수록 예외 없이 작다} --- 쇠퇴비는 3분의 1인데 날짜 상관이 세 배다.

\begin{figure}[t]
\centering
\includegraphics[width=\columnwidth]{figs/clock.pdf}
\caption{왼쪽 --- 도메인마다 모형과 달력. 오른쪽 --- 쇠퇴비와 시계 몫이
따로 논다.}
\end{figure}

\section{시계 몫을 박았다}

쓸 눈금은 \textbf{달력 $\div$ 모형}이다.

\[ \text{시계 몫} = \frac{\rho(-\text{시작일},\ y)}{\rho(\hat{y},\ y)} \]

1 을 넘으면 날짜 한 줄이 축 스물셋보다 낫다는 뜻이다.
\texttt{rank.clock} 으로 박았다 --- 순위에 안 쓰고 도메인 수를 읽는 눈금이다
(\texttt{rank.spread} · \texttt{rank.by\_kind} 와 같은 자리).

\begin{center}\small
\begin{tabular}{lrrr}
\hline
도메인 & 모형 & 달력 & 시계 몫 \\
\hline
\textbf{도서} & $+$0.369 & $+$0.561 & \textbf{1.52} \\
웹툰 & $+$0.453 & $+$0.350 & 0.77 \\
모바일 & $+$0.535 & $+$0.198 & 0.37 \\
게임 & $+$0.622 & $+$0.182 & 0.29 \\
애니 & $+$0.505 & $+$0.144 & 0.28 \\
세계애니 & $+$0.534 & $+$0.099 & 0.19 \\
펀딩 & $+$0.267 & $-$0.012 & $-$0.04 \\
팝업 & $+$0.361 & $-$0.113 & $-$0.31 \\
\hline
\end{tabular}
\end{center}

\textbf{여덟 중 도서 하나만 1 을 넘는다.} 웹툰이 0.77 로 둘째인데 1 아래라
모형이 이긴다 --- 노트 267에서 코호트 통제에도 웹툰 모형이 0.441 $\to$
0.426 으로 거의 안 움직인 것과 맞는다. 팝업 · 펀딩은 음수다: 오래 된 것이
오히려 낮다(팝업은 코로나 이전 방문자가 적었고, 펀딩은 캠페인이 끝나면
후원자가 더 안 는다).

\section{두 가지를 곁가지로 확인했다}

\paragraph{게임 잘림 처리는 옳다} \texttt{y\_w30} 은 출시 30일이 안 찬
작품에서 의미가 없다. 레코드에 \texttt{y\_w30\_truncated} 가 있고
\texttt{fixaxes.\_labels} 가 그 플래그가 서면 라벨을 결측으로 둔다. 확인하니
\textbf{43건 전부} 판에서 NaN 이다. 노트 210의 고침이 온전하다.

\paragraph{그 43건이 내 코드를 물었다} \texttt{rank.clock} 을 처음 돌렸을 때
\textbf{게임이 표에서 빠졌다.} 원인은 내가 ``유보 행 수''와 ``라벨이 있는
행 수''를 같은 것으로 보고 예측 벡터 길이와 견줬기 때문이다 ---
\texttt{predictions} 는 유보 행을 \emph{전부} 돌려주고(라벨 NaN 포함) 게임은
그 차이가 정확히 43 이었다. 조용히 한 도메인이 사라졌고 예외도 안 났다.

노트 264(가드가 마스크를 안 봤다) · 노트 266(후보 검사가 결측 행을 넣었다)에
이어 \textbf{세 번째로 같은 자리에서 넘어졌다} --- 새로 쓴 재는 코드가
``어느 행이 진짜 있는 행인가''를 잘못 잡았다. 셋 다 조용히 틀렸고 셋 다
숫자가 이상해서 알아챘다.

\paragraph{남은 것}
\begin{center}\small
\begin{tabular}{ll}
\hline
갈래 & 왜 \\
\hline
웹툰 0.77 & 둘째로 높다. 몫 27\% 라 판에 제일 크게 실린다 \\
행 잡기 & 세 번 틀렸다. 공통 도구로 뺄까 \\
만화 · 아이돌 & 표본 부족으로 이 검사를 못 한다 \\
\texttt{SOURCE} & 34/45 쌍 \\
\hline
\end{tabular}
\end{center}

\paragraph{규약} 라벨 오염을 잴 때 \textbf{쇠퇴의 크기가 아니라 순위가 날짜를
따르는 정도}를 본다. 게임은 34배 쇠퇴해도 멀쩡하고 도서는 10배에 무너진다.
\texttt{rank.clock} 을 도메인 수와 함께 읽는다.

\end{document}
